% tab:previous
\begin{table}[htbp]
\centering\small
\caption{Latvian GDP per capita: the new series and earlier estimates, selected years}
\label{tab:previous}
\begin{threeparttable}
\begin{tabular}{lrrrr}
\toprule
Year & This paper & \multicolumn{3}{c}{Earlier estimates (difference from this paper, \%)} \\
\cmidrule(lr){3-5}
& (\intlnineteen{}) & Ros\'{e}s and Wolf (2010) & Norkus and Markevi\v{c}i\={u}t\.{e} (2021) & MPD 2023 \\
\midrule
1913 & 2,361 & --- & 2,361 (+0.0) & --- \\
1922 & 1,654 & 1,929 (+16.6) & 1,847 (+11.7) & --- \\
1929 & 2,351 & 2,798 (+19.0) & 2,347 ($-$0.2) & --- \\
1938 & 3,099 & 4,048 (+30.6) & 2,836 ($-$8.5) & --- \\
1973 & 8,523 & --- & --- & 7,846 ($-$7.9) \\
1980 & 8,867 & --- & --- & 8,619 ($-$2.8) \\
1985 & 9,284 & --- & --- & 8,954 ($-$3.6) \\
1989 & 8,772 & --- & --- & 9,825 (+12.0) \\
1990 & 7,965 & --- & --- & 9,916 (+24.5) \\
\bottomrule
\end{tabular}
\begin{tablenotes}[para,flushleft]
\footnotesize
\textit{Notes:} MPD figures are published in \intltwenty{} and converted at 1.594. The 1913 figure of this paper is that of Norkus and Markevi\v{c}i\={u}t\.{e}, taken as published. \textit{Sources:} \textcite[190]{roseswolf2010}; \textcite[603]{norkusmark2021}; \textcite{mpd2023}; this paper (online appendix).
\end{tablenotes}
\end{threeparttable}
\end{table}
